\section{Data and Methodology}

\subsection{Data and Sample Selection}
Our dataset spans a ten-year period from January 1, 2015, to December 31, 2024. It consists of two main components:
\begin{enumerate}
    \item \textbf{Index Price Data:} Daily Open, High, Low, and Close prices for the VN-Index (HoSE benchmark) were obtained from Quote history files. The daily log return is computed as $r_t = \ln(Close_t / Close_{t-1})$.
    \item \textbf{News Corpus:} A web-scraped database of financial articles from the CafeF sitemap. After hash-based deduplication and length filtering ($body\_len \ge 100$ characters), the corpus contains $86,179$ unique articles.
\end{enumerate}

\subsection{Trading-Day Alignment and Imputation}
Because financial markets operate on active trading sessions whereas sitemaps publish news continuously (including nights, weekends, and holidays), proper temporal alignment is critical. We define a strict trading-day alignment rule based on the HoSE session closure:
\begin{itemize}
    \item Articles published on a trading day $t$ before or at \textbf{14:45} (session close) are aligned to the current trading day $t$.
    \item Articles published on a trading day $t$ after \textbf{14:45}, on weekends, or during public holidays (e.g., Lunar New Year) are aligned to the next active trading day $t+1$.
\end{itemize}

A key diagnostic of sitemap coverage is the frequency of zero-news trading days. The daily panel contains $2,498$ trading-day rows, and after the one-step-ahead shift there are $2,497$ forecastable volatility targets. Only $4$ days ($0.16\%$) have zero aligned news articles, so missing daily sentiment features are handled with zero-imputation:
$$\text{n\_articles}_t = 0, \quad \text{mean\_sentiment}_t = 0.0, \quad \text{sentiment\_std}_t = 0.0$$

\subsection{Sentiment Construction}
The CafeF sentiment pipeline first runs PhoBERT classifier inference at the article level and writes \texttt{sentiment\_score} as the difference between positive and negative class probabilities. The article output also records \texttt{sentiment\_label} and class probabilities for validation. Daily features are then aggregated as mean sentiment, sentiment dispersion, sentiment volume, label shares, net sentiment, a five-day sentiment surprise measure, and category-specific macro and market sentiment.

The sentiment score is a proxy rather than a hand-labeled gold standard, so the interpretation in later sections is descriptive rather than causal.

\subsection{Volatility target definitions}
To evaluate forecast accuracy, we construct two daily volatility proxies:
\begin{enumerate}
    \item \textbf{Parkinson Volatility:} A high-low range estimator:
    $$\sigma_{P, t} = \sqrt{\frac{\left(\ln\left(H_t / L_t\right)\right)^2}{4 \ln 2}}$$
    where $H_t$ and $L_t$ are the high and low prices of index on day $t$.
    \item \textbf{Garman-Klass Volatility:} An estimator incorporating opening and closing gaps:
    $$\sigma_{GK, t} = \sqrt{0.5 \left(\ln\frac{H_t}{L_t}\right)^2 - (2\ln 2 - 1)\left(\ln\frac{C_t}{O_t}\right)^2}$$
    where $C_t$ and $O_t$ are the close and open prices.
\end{enumerate}
The target variable is the lead volatility proxy: $y^{\text{vol}}_{t+1} = \sigma_{t+1}$.

\subsection{Econometric \& Neural Network Specifications}

\subsubsection{Baseline AR(1)-GARCH(1,1)}
Returns are scaled by 100 ($y_t = r_t \times 100$). The mean equation is modeled as an autoregressive process of order 1 to capture short-term return autocorrelation:
$$y_t = c + \phi y_{t-1} + e_t$$
where $c$ is a constant intercept, $\phi$ is the autoregressive parameter, and $e_t \sim N(0, \sigma_t^2)$ is the mean-equation residual. The conditional variance of the residuals is modeled recursively:
$$\sigma_t^2 = \omega + \alpha e_{t-1}^2 + \beta \sigma_{t-1}^2$$
subject to constraints: $\omega > 0$, $\alpha \ge 0$, $\beta \ge 0$, $\alpha + \beta < 1.0$, and $|\phi| < 1.0$ for covariance stationarity.

\subsubsection{AR(1)-GARCH-X(1,1)}
For comparison, the GARCH-X model integrates the exogenous sentiment variable $X_{t-1}$ directly into the variance equation:
$$\sigma_t^2 = \omega + \alpha e_{t-1}^2 + \beta \sigma_{t-1}^2 + \gamma X_{t-1}$$
where $e_{t-1}$ is the lagged residual from the AR(1) mean equation and $X_{t-1}$ is the lagged daily aggregate \texttt{mean\_sentiment}.

\subsubsection{Hybrid GARCH + Sentiment LSTM}
The hybrid model uses a two-stage approach. In Stage 1, baseline AR(1)-GARCH parameters are estimated, and the 1-step-ahead GARCH volatility forecast $\sigma_{t, \text{GARCH}}$ (standard deviation) is computed. The conditional residual target for day $t$ is:
$$\text{hybrid\_residual\_target}_t = \sigma_{t, \text{realized}} - \sigma_{t, \text{GARCH}}$$
In Stage 2, an LSTM network is trained to predict this residual. The input sequence is a rolling window of length $L=15$ days containing 17 features (including returns, GARCH standardized residuals $\epsilon_t$, daily article counts, daily mean sentiment, net sentiment, macro/market category sentiment shares, and GARCH conditional variance $\sigma^2_t$). Crucially, the LSTM inputs utilize the conditional variance $\sigma^2_t$ (via the \texttt{garch\_forecast\_var} feature) rather than the standard deviation ($\sigma_t$) to represent historical baseline risk.
The final hybrid volatility forecast is:
$$\sigma_{t, \text{hybrid}} = \sigma_{t, \text{GARCH}} + \hat{u}_{t, \text{LSTM}}$$
where $\hat{u}_{t, \text{LSTM}}$ is the residual predicted by the neural network.

The current implementation uses a chronological train/test split (with the validation set size set to 0 to maximize training samples) and treats the held-out window as the evaluation sample.

